\paragraph{} Visually rich business documents such as contracts, invoices, reports, and scanned forms remain difficult for automated systems because their information is distributed across text, layout, tables, page structure, and domain-specific field conventions. Recent document AI models improve extraction and document understanding, while Retrieval-Augmented Generation (RAG) improves access to external knowledge. However, practical document question answering still suffers from three research problems: extraction uncertainty propagates into retrieval, retrieved evidence may be only partially relevant, and language models can generate unsupported or contradictory answers even when retrieval is used.

\paragraph{} This thesis investigates the research question of how uncertainty-aware document extraction and evidence-graded Agentic RAG can improve trustworthy document intelligence for semi-structured and visually rich business documents. The implemented system is used as a research prototype to study an end-to-end pipeline that combines OCR, layout-aware document analysis, structured field extraction, validation, semantic retrieval, evidence grading, and grounded generation. Unlike a direct OCR-to-RAG baseline, the proposed workflow introduces query planning, structured-field shortcutting, candidate evidence retrieval, evidence scoring, confidence estimation, and traceable answer generation.

\paragraph{} The central hypothesis is that document question answering becomes more reliable when the system explicitly separates four decisions: whether a question can be answered from structured extractions, which chunks should be retrieved, which retrieved evidence is sufficiently relevant, and whether the final answer is grounded enough to return. This framing is aligned with recent research on practical document understanding benchmarks, RAG evaluation, groundedness, hallucination detection, and self-reflective retrieval-augmented generation.

\paragraph{} A preliminary prototype evaluation reports a field extraction F1 score of 0.8, Search Precision@5 of 0.25, Search Mean Reciprocal Rank of 1.0, estimated throughput of 29.77 documents per hour, and a review rate of 0.152. These early results are treated as pilot evidence rather than final claims. The thesis further defines a research-oriented evaluation protocol comparing multiple baselines: OCR-only extraction, OCR plus direct RAG, structured extraction plus direct RAG, and the proposed evidence-graded Agentic RAG workflow. The proposed metrics include extraction F1, Precision@k, MRR, answer faithfulness, context precision, context recall, unsupported-claim rate, latency, and human-review load.

\paragraph{} The work contributes a reproducible prototype, a research framework for uncertainty-aware document QA, a baseline comparison plan, and a PhD-ready research agenda for trustworthy multimodal document intelligence. Future work will expand the benchmark dataset, add labelled question-answer pairs, introduce stronger table and layout modelling, fine-tune document understanding models, and evaluate faithfulness and hallucination behaviour at scale.
